\section{应用与推广}\label{sec:applications}

\subsection{算术几何的标准技术}

怀尔斯证明发展的技术已成为当代数论的日常工具：

\begin{itemize}
    \item \textbf{模性提升（$R=T$ 定理）}：如今证明“某伽罗瓦表示来自模形式”的标准模板，应用于 Serre 水平约化猜想的证明（Khare--Wintenberger 2009\cite{khare2009serre}）、Sato--Tate 猜想的证明（Clozel--Harris--Shepherd-Barron--Taylor 等 2008\cite{clozel2008sato}）、以及自守表示的若干提升问题；
    \item \textbf{欧拉系}：Kato 欧拉系与 Selmer 群的结构定理被推广到 BSD 猜想、主猜想的研究；
    \item \textbf{Hecke 代数的完全交性质}：完全交判别法与 Gorenstein 条件的运用已成为 Hecke 代数计算的标准假设。
\end{itemize}

\subsection{朗兰兹纲领}

模性定理是朗兰兹纲领在 $\operatorname{GL}_2/\mathbb{Q}$ 情形的核心实例：伽罗瓦表示 $\leftrightarrow$ 自守表示的对应。怀尔斯的证明提供了“数域上二维伽罗瓦表示可以自守化”的第一个完整范例，此后的推广（全实域上的 $\operatorname{GL}_2$、$\operatorname{GL}_n$ 的潜在自守性等）都建立在同一技术框架上。费马大定理因此可以视为朗兰兹纲领威力的第一个重大展示。

\subsection{椭圆曲线计数与密码学的联系}

椭圆曲线密码学（ECC）的安全性依赖椭圆曲线离散对数问题的困难性。模性定理保证了所有 $\mathbb{Q}$ 上椭圆曲线的 $L$ 函数具有良好的解析性质，这在曲线选取（如 CM 方法、安全素数阶群检验）与离散对数攻击分析（如 anomalous 曲线的排除）中提供理论基础。费马大定理本身对密码学没有直接应用，但其证明所依赖的基础设施（Hasse 界的推广、点 counting 的有效算法）渗透在 ECC 的工程实践中。

\subsection{丢番图方程的统一视角}

FLT 的解决确立了“把丢番图方程翻译为伽罗瓦表示与模形式”的范式。该范式现在被称为 modular method（模方法），已系统应用于：

\begin{itemize}
    \item 广义费马方程 $A x^p + B y^p = C z^p$ 的特定族（Bennett--Skinner、Dahmen 等的 modular method 计算）；
    \item 完全立方、双二次与 Lebesgue--Nagell 型方程；
    \item 椭圆曲线上的积分点问题与 Frey--Hellegouarch 构造的系统化。
\end{itemize}

\subsection{数学文化的影响}

从1993年 BBC 的直播讲座到西蒙·辛格的畅销书，费马大定理的解决过程极大提升了公众对纯数学的关注。它也是“数学合作的长期主义”的象征：三百五十七年间，欧拉、库默尔、弗雷、塞尔、里贝特、马祖尔等人的工作层层接力，最终由怀尔斯收束——单个天才的神话背后，是整个学科几代人的基础设施。
